 S_* =\{p:p\le76176\},\qquad
 \mathcal D=\{0,1,2401,-343/32\},
\]
\[
 H_*=2(4608^{165888}+6^{100})
          \left(\prod_{p\le76176}p\right)^{24}+128.
\]
For $a\in\Q\setminus\mathcal D$, let $M_5(a)$ be the product of
the primes outside $S_*$ at which the point $(a:1)$ has positive odd
contact order with one of $\infty,0,1,2401,-343/32$. Explicitly, choose
local homogeneous coordinates $(u:v)$ and subtract
$\min(v_p(u),v_p(v))$ from the valuations of $v$ and of
$u-\delta v$ for $\delta\in\mathcal D$.

\begin{theorem}[The five-contact height estimate]\label{hyp:o8-height}
For every $a\in\Q\setminus\mathcal D$,
\[
 h(a)\le H_*M_5(a)^{24}.
\]
The displayed constant satisfies $\log_2H_*<4783227$.
\end{theorem}
\begin{proof}
Write $s=a$, $M(s)=M_5(a)$, $H=H_*$,
$B_*=\prod_{p\le76176}p$, and
$\mathcal C=\{\infty,0,1,2401,-343/32\}$.
We prove the geometric comparison, the non-CM estimate through the
modular level, and the CM estimate over the same fiber field.
All heights use the original stable Faltings normalization.
\paragraph{Recovering the contact line from published quotients}
Let $B/\mathbb Q$ be the indefinite quaternion algebra of discriminant $210$,
$\mathcal O$ a maximal order, $X$ its maximal-level Shimura curve, and $W$
its Atkin--Lehner group. The standard genus formula gives
\[
 g(X)=1+\frac{\varphi(210)}{12}=5,
\]
because the elliptic-point terms vanish: $5$ splits in $\mathbb Q(i)$ and
$7$ splits in $\mathbb Q(\sqrt{-3})$, so neither quadratic field embeds in
$B$. The group $W$ is $(\mathbb Z/2\mathbb Z)^4$ with involutions $w_m$
indexed by square classes of divisors of $210$.

The published models of \cite[Theorem 4.2 and Table 2, pp.~938--939]{GonzalezRotger2006} have the following coefficient rows
$(a,b,c,d,e)$ for $y^2=ax^4+bx^3+cx^2+dx+e$:
\[
\begin{array}{c|rrrrr}
 m&a&b&c&d&e\\\hline
30&-43&-686&-2915&-1372&-172\\
42&-43&600&-986&-600&-43\\
70&-43&256&-130&-768&-387\\
105&-43&2&505&12&-1548
\end{array}
\]
The cited theorem identifies each with $X/\langle w_m\rangle$.
The genus formula and the Atkin--Lehner group law used here also appear
in the same published paper, Proposition 5.2 and the paragraph following
Proposition 5.3, pp.~940--941.

Here is a deduction of all the required branch information, rather than
an additional assumption about it. Each of these four genus-one quotients
has eight fixed points upstairs, by Riemann--Hurwitz. A stabilizer for the
$W$ action on a characteristic-zero smooth curve is cyclic, and therefore
has order at most two. Each nonempty fixed set of an involution is a union
of $W$-orbits of size eight. Riemann--Hurwitz for $X\to X/W$ now gives
\[
 8=16(2g(X/W)-2)+\sum_{1\ne w\in W}\#\operatorname{Fix}(w).
\]
The four known fixed sets already contribute $32$. Thus $g(X/W)=0$,
the total is $40$, and there is exactly one further orbit of eight fixed
points. Inertia generators around the branch points have product one.
The four classes $30,42,70,105$ are independent in $W$, and their product
is the class $210$. Hence there are exactly five branch points, each of
index two, with distinct labels $30,42,70,105,210$. Since the involutions
are defined over $\mathbb Q$ and each label occurs only once, each branch
point is rational. Three of them trivialize the genus-zero quotient.

Normalize $b_{210}=\infty$, $b_{105}=0$, $b_{70}=1$, and write
$b_{42}=\lambda$, $b_{30}=\mu$. On a genus-one quotient the residual
group $(\mathbb Z/2)^3$ has translation subgroup of order four, necessarily
the full two-torsion group after choosing an origin. Its quotient has the
same Jacobian $j$-invariant: it is multiplication by two over the algebraic
closure. The remaining double cover of $X/W$ branches at the four points
other than $b_m$. Therefore the four quartic invariants determine the four
Legendre invariants of these four-point subsets.

For a quartic put
\[
 I=12ae-3bd+c^2,\quad
 J_q=72ace+9bcd-27ad^2-27b^2e-2c^3,\quad
 j=\frac{6912I^3}{4I^3-J_q^2}.
\]
The exact computation gives the Legendre parameters
$1/2401,32/375,1/225,5/1029$, in row order. Enumerating the six cross-ratio
values of each of the first two parameters and testing the last two
invariants leaves exactly
\[
 \lambda=2401,\qquad\mu=-343/32.
\]
The script \texttt{check\_height.py} reproduces this rational computation.
Every nonzero determinant of the five primitive contact forms has prime
support in $\{2,3,5,7\}$.

\paragraph{Integral geometry and the degree 576 field}
Choose $\mu_0\in\mathcal O$ with $\mu_0^2=-210$. The embedding criterion
allows this choice in a maximal order; maximal orders are conjugate here.
Use the maximal moduli stack $\mathscr X$ and the entire principal-level
three scheme $Y$, including its geometric components. Yuan's integral
moduli theorem \cite[Theorem 1.1 and Sections 2.1--2.2]{Yuan2024} makes $\mathscr X$ regular, proper and semistable over
$\mathbb Z$, with a universal abelian surface, and smooth away from $210$.
Its generic stabilizer is precisely $\{\pm1\}$. Principal level three is
fine and gives a universal full three-torsion marking.
